\section*{The Propers: Detailed Commentary}

\subsection*{Trust selected from inside lament\enspace\elementref{Entrance; Collect}}

Psalm 54's confession of sustaining help is not cheerful detachment from danger. The full psalm knows terror, violent city life, and betrayed friendship. Appointed at the Entrance, its trust is costly: God upholds the speaker who cannot secure the world. The accompanying resolve to offer and praise makes worship response rather than leverage.

Augustine's exposition follows the Vulgate/LXX Psalm 53 and locates the help “within,” in the conscience formed toward God rather than in visible success. He then reads the promised voluntary sacrifice as praise offered because God is good, not as flattery exchanged for prosperity. His reading bears on the adapted Entrance: divine support is present before danger is resolved, and praise is not a technique for forcing the desired outcome.

The Collect asks God to look favorably upon servants and multiply gifts of grace. Hope, faith, and charity are not interchangeable moods; they are theological virtues whose object and source are God. Vigilance over the commands therefore follows gift. The prayer avoids both self-generated virtue and grace imagined as exemption from obedience.

\subsection*{Clemency from strength, not from denial\enspace\elementref{First Reading; Responsorial Psalm}}

Wisdom 12 addresses a classic temptation: to imagine mercy and power as competitors. The text reverses that assumption. Because God's mastery is complete, he need not prove himself through anxious severity. His might is the source of justice, and his government can be lenient while remaining judgment. The result is pedagogical: the just learn kindness, and sinners receive grounds for hope in repentance.

Augustine quotes the core Wisdom 12:15--18 sequence in \emph{Eighty-three Questions} 53.3 while asking how God's tranquil judgment works through ministering agents. He treats the verses directly rather than commenting continuously on Wisdom's literary unit: divine action lacks perturbation even when created ministers experience and execute judgment in time. In his expositions of Psalms 2 and 6 he returns to verse 18 to distinguish divine judgment from the passions that disfigure human anger.

Gregory the Great likewise quotes Wisdom 12:18 in \emph{Moralia in Job} V.78. His moral reuse joins tranquil judgment to restraint of anger: the ruler who tempers severity with mildness returns toward the likeness of the Creator. Gregory does not dissolve judgment into indulgence, nor does he expound Wisdom's Alexandrian literary setting. Augustine and Gregory together witness to a Latin reception in which divine mastery is untroubled and therefore trains creaturely authority away from reactive domination.

Salvian's \emph{On the Governance of God} II quotes Wisdom 12:15 and 18 while arguing that universal divine government gives human beings dignity and makes their moral answer consequential. Aquinas later uses verse 18 in his \emph{Exposition on the Lord's Prayer}: rational creatures are governed not merely as things moved from outside but as children capable of knowing and loving the Father's will. These Latin receptions move in different genres---apologetic Providence and catechesis on prayer---yet both protect the appointed conclusion that divine patience is pedagogical. Hope of repentance is not mere delay; it addresses persons as responsible creatures who can be taught to act justly and kindly.

The sequence Augustine--Gregory--Salvian--Aquinas also exposes a limit. “Judge with clemency” does not pass unchanged from divine to human power, which must reckon with creaturely ignorance, office, law, and the duty to protect. The Fathers and Doctor offer analogies of unperturbed judgment and filial governance, not a warrant for officials to claim God's mastery of every heart and consequence.

The omitted verses matter. Lectionary 106 moves from the assertion that no higher tribunal audits God directly to the power--justice relation, leaving 12:14--15 outside proclamation. That wider argument informs the reading, but the assembly does not hear every step.

Psalm 86 transforms teaching into prayer. The servant appeals to a God who is good, forgiving, merciful, gracious, patient, and faithful, while confessing that all nations will worship him. Mercy reaches the servant and all nations. The response asks the Judge described by Wisdom to act according to his revealed character.

\Needspace{6\baselineskip}
Augustine's exposition of the Vulgate/LXX Psalm 85 opens with a Christological account of prayer: Christ prays for us as priest, prays in us as head of the Body, and is prayed to as God. Later, when the psalm asks for an undivided heart, Augustine diagnoses the mind that wanders even while the mouth prays. That reception gives the responsorial action moral depth. Repeating that God is good and forgiving is not merely reporting an attribute; it asks that distracted desire be gathered into the mercy it invokes.

\subsection*{Weak prayer genuinely aided\enspace\elementref{Second Reading}}

Romans 8:26--27 continues a sequence of groaning: creation, believers, and now inexpressible prayer. Weakness does not mean that every desire is pure or that discernment is unnecessary. It means finite, suffering persons cannot fully know the fitting shape of petition within God's whole design.

The Spirit “comes to the aid” rather than rendering the person irrelevant. Prayer remains a human act made possible and ordered from within by divine agency. The Father who searches hearts knows the Spirit's intention; the verse depicts Trinitarian communion, not a contest in which the Spirit persuades an unwilling God.

Chrysostom's \emph{Homily 14 on Romans} reads the “groanings” in relation to the Spirit-given prayer of the Church and refuses to make spiritual help an excuse for passivity. Augustine returns to the same verses in \emph{Letter 130} 15 and 28--29. He calls the disciple's condition a learned ignorance: believers know the end they seek but do not know which temporal form will serve it, so prayer may be carried more by tears and groaning than by many words. In \emph{On the Soul and Its Origin} IV.13 he uses Paul's unanswered request concerning the thorn as a concrete case: a sincere temporal petition can miss what is fitting, while the Spirit forms the saint's deeper intercession.

These witnesses converge yet differ exegetically. Chrysostom attends closely to charismatic, ecclesial prayer; Augustine stresses the transformed desire of the praying saint. Neither imagines the Spirit uttering information unknown to the Father. Both preserve human weakness, divine initiative, and the real act of prayer.

Aquinas's account of created causes within Providence, read with the Catechism's treatment of secondary causality, supplies a doctrinal boundary. Prayer can be a real secondary cause within Providence; the Spirit's intercession keeps petition ordered toward the true good even when a requested temporal form is withheld. This does not make words, liturgy, prudent means, or persevering request unnecessary. Silent anguish and articulated prayer can both be taken into the Spirit's aid.

\subsection*{The short and long Gospel paths\enspace\elementref{Acclamation; long Gospel; short Gospel}}

The acclamation adapts Matthew 11:25 and welcomes the Kingdom discourse through receptive littleness. It does not say that mature learning is evil; it blesses revelation as gift rather than possession. Its narrative source is earlier than Matthew 13 and lies outside the house where Jesus later explains the weeds.

Chrysostom's \emph{Homily 38 on Matthew} directly treats this thanksgiving. “Wise and prudent” names self-sufficient resistance, not the disciplines of learning, while “little ones” receive because they do not pretend to stand above gift. The adapted acclamation therefore asks for teachability before a parable that will frustrate premature certainty; it gives no weapon against scholarship, careful safeguarding, or the Church's duty to discern.

The short Gospel tells a complete parable. Good seed is sown; an enemy introduces weeds; servants discover the mixture and propose removal; the master forbids it because wheat could be uprooted; both remain until harvest, when distinct destinations are assigned. Patience protects the good while preserving separation at the proper time.

Chrysostom's \emph{Homily 46 on Matthew} distinguishes this parable from the preceding sower. Here the seed has taken root and a counter-sowing imitates it; evil works by mixture and counterfeit rather than creation. He reads the prohibition against uprooting as a ban on killing persons for error and notes the danger that wheat will fall with the weeds. His homily still permits checking falsehood and defending truth. It therefore supports both patient hope and safeguarding, not an undifferentiated policy of institutional inaction.

The long form adds two different growth images. Mustard begins almost invisibly and becomes capacious; leaven works hidden within a large quantity of dough. Matthew then explains Jesus' parabolic speech with an internal Psalm citation and moves from crowd to house. Jesus identifies sower, field, seed, enemy, harvest, and harvesters, and culminates in the removal of evil and the righteous shining in the Father's kingdom.

\Needspace{4\baselineskip}
The same Chrysostom homily reads the apostles in the mustard and leaven: few and outwardly weak, they are entrusted with a word capable of reaching the whole. That missionary reception belongs only to the long branch. It also keeps “hidden growth” from becoming an excuse for doing nothing; the leaven is hidden because it is working.

Origen's surviving Greek \emph{Commentary on Matthew}, Book X, begins its sustained treatment at the private explanation in 13:36. He notices the spatial movement from crowd to house and reads sleeping as failure to keep watch while hostile teaching is sown among good conceptions. Unlike a modern botanical decoding, his account is a spiritual-intellectual diagnosis of counterfeit doctrine. His reading is branch-specific: a short-form congregation does not hear the house, enemy, devil, or interpretive key.

Chrysostom's \emph{Homily 47} follows the explanation to judgment. The Son owns the field, sends the angels, removes causes of sin, and brings the righteous into the Father's kingdom. Patience is thus bounded by eschatology: delayed separation is neither denial that evil is evil nor assurance that mixture is permanent.

Those additions materially change what is heard. A short-form congregation does not hear the devil, angels, fiery furnace, or shining righteous named, though canonical context still bears on the parable. Conversely, the short form's delayed harvest does not suppress the long form's final judgment within Christian doctrine.

\begin{studybox}{Patience is not a safeguarding policy}
The parable denies servants the Master's infallible final separation; it does not forbid removing a dangerous person from access, reporting abuse, applying proportionate discipline, protecting victims, correcting doctrine, treating disease, or resisting injustice. Such actions are governed by other moral, canonical, civil, and pastoral norms. “Do not uproot the wheat” is itself a protection command, not permission to expose wheat to preventable harm.
\end{studybox}

Nor is “children of the evil one” a label that authorizes contempt. Jesus describes the eschatological meaning of the parable. Living persons remain objects of evangelization and possible repentance; Wisdom's appointed conclusion explicitly gives sinners hope of turning back. God alone judges hearts without error.

\subsection*{One sacrifice and common salvation\enspace\elementref{Prayer over Offerings}}

The Prayer over the Offerings makes a strong historical-theological claim in compressed Latin: the diversity of sacrifices under the Law is sanctioned in the perfection of one sacrifice. It then asks God to receive servants' gifts and order them to the salvation of all. Individual offering is therefore neither private possession nor another sacrifice alongside Christ's. It is gathered into the one sacramental action and directed toward common good.

\subsection*{Memorial food or a repentant opening\enspace\elementref{Communion A; Communion B; Prayer after Communion}}

Psalm 111's Communion alternative joins remembered wonders, mercy, and food. Its full acrostic praise also names covenant, redemption, justice, and reverence. Eucharistic memorial is more than mental recall, but the antiphon's psalmic source rightly situates reception within grateful remembrance of God's works.

Augustine's exposition of the Vulgate/LXX Psalm 110 reads God's memorial and food together. He receives the food given to those who fear God as bread from heaven and a pledge of the inheritance, while keeping reverence and covenant in view. His Eucharistic reading goes beyond the psalm's original historical referent and bears on branch A alone, yet it helps explain why memorial at Communion cannot be reduced to recollection of absent events.

Alternative B takes Revelation 3:20 from the address to Laodicea. The familiar knock is not first a sentimental image of an isolated individual. The risen Christ rebukes a self-satisfied church, commands zeal and repentance, then promises shared table fellowship. At Communion this branch joins invitation to conversion; sacramental approach cannot be detached from hearing the one who knocks.

Victorinus of Pettau's surviving \emph{Commentary on the Apocalypse} treats the Laodicean rebuke but does not furnish a separate exposition of verse 20 in the received English text. It can corroborate the corporate setting and summons from lukewarmness to repentance, not a detailed Eucharistic interpretation of the knock. Liturgical appointment makes the verse apt for Communion, while Victorinus keeps its admonitory context in view.

\Needspace{6\baselineskip}
Ambrose receives the verse directly in \emph{Exposition of the Christian Faith} IV.19--23. He hears Christ standing at the door of the soul and joins the image to the beloved's knock in the Song and the Psalm's gates lifted for the King of glory. Jerome, \emph{Letter 22} 26, and Cassian, \emph{Institutes} V.19--20, likewise use the verse to press prompt interior response rather than delayed conversion. These witnesses receive the knock as Christ's summons to open, enter fellowship, and overcome the habits that bar the door.

The two antiphons remain distinct textual paths: one proceeds through remembered wonders and provision, the other through rebuke, repentance, and promised fellowship. The Prayer after Communion asks that the people imbued with heavenly mysteries pass from oldness into new life. “New” does not mean contempt for creation or memory; it names Paschal transformation of persons and conduct.

\section*{The Propers: Notable and Quotable}

{\small
\begin{itemize}
\item \textbf{“You might uproot the wheat along with them” --- a charter for unpersecuted consciences.\enspace\elementref{both Gospel forms}} Roger Williams's \emph{Bloudy Tenent of Persecution} (1644), chs.~XVIII--XXVIII, expressly argues from Matthew 13:30, 38. In ch.~XXVI the endangered wheat becomes a civil-peace argument: punish civil offenses, he says, but do not uproot worship and conscience. The harvest image has become political reasoning, not a farming illustration.

\item \textbf{“Let them grow together until harvest” --- an objection to licensing ideas.\enspace\elementref{both Gospel forms}} In \emph{Areopagitica} (1644), John Milton says mortals cannot “sever the wheat from the tares”; that belongs to angels at the end. He turns delayed sorting into a case against compulsory conformity and prior licensing: human censors may destroy truth while trying to remove error. His own toleration still had sharp historical exclusions.

\item \textbf{“Wailing and grinding of teeth” --- with dental provision guaranteed.\enspace\elementref{long Gospel}} Irvin S.~Cobb's 1923 joke collection, no.~358, has an elderly hearer interrupt a damnation sermon because she has no teeth. The preacher's answer---“Teeth will be provided!”---literalizes a figure of eschatological anguish into a brisk supply problem. Cobb's wording blends the familiar formulas from several Matthean passages; dependence on Matthew 13:42 alone cannot be shown.

\item \textbf{“I stand at the door and knock” --- a door with no outside handle.\enspace\elementref{Communion B}} William Holman Hunt's \emph{The Light of the World} (1851--1853) paints Revelation 3:20 as Christ before an overgrown door that can be opened only from within. The visual addition shifts a rebuke addressed to a complacent church into an image of the individual heart or mind and its agency; the Museum of Oxford records that the picture became one of its age's most widely reproduced images.

\end{itemize}
}

\section*{The Propers: Interpretive Possibilities}

\noindent\emph{The following exploratory proposals were developed in an AI-assisted editorial process. They are attributed to none of the cited authorities and claim no historical compositional intent. Each states its anchors and limits.}

\Needspace{10\baselineskip}
\noindent\textbf{Learned ignorance can restrain the hand without paralyzing it.\enspace\elementref{First Reading; Second Reading; Gospel}}
Wisdom's ruler judges clemently toward repentance; Matthew's servants see weeds but cannot uproot them safely; Romans joins ignorance in prayer to divine knowledge and intercession. Together they distinguish recognizing evil from controlling every heart or outcome. Separate reading misses prayer as action within that gap. Its fruit is judgment that protects, reports, corrects, and prays without claiming final knowledge. Augustine's “learned ignorance” and official patient-hope preaching are near analogues. Romans is semi-continuous; uncertainty never excuses evidence, competent authority, safeguarding, or action against imminent harm.

\Needspace{10\baselineskip}
\noindent\textbf{Hiddenness is not a moral verdict.\enspace\elementref{long Gospel; Acclamation}}
Only the long Gospel juxtaposes nocturnal counterfeit seed, a minute seed becoming shelter, and hidden leaven transforming meal; the acclamation adds revelation to little ones. Hiddenness can conceal evil or mark the Kingdom, so scale, secrecy, and speed cannot identify moral source. Separate reading misses this opposed hidden corruption and fecundity. Its fruit is discernment by truth, fruit, relation to Christ, and accountable process rather than virality. Chrysostom's sequence is a near analogue but lacks the diagnostic. Long-form only: good and evil remain distinguishable, and doctrine, morals, evidence, accountability, and safeguarding govern.

\Needspace{10\baselineskip}
\noindent\textbf{The one sacrifice interrupts the manufacture of scapegoats.\enspace\elementref{First Reading; Gospel; Prayer over Offerings}}
Wisdom's strength gives hope of repentance; the householder protects wheat from uprooting; the offering prayer gathers many sacrifices into one perfect sacrifice for all. Together they forbid purification that offers expendable persons to communal fear or unity. Separate reading misses Christ's sacrifice judging scapegoating as false liturgy. Its fruit is a test for polarization, self-protection, and punitive spectacle that retains accountability. The exact conjunction was not located in the checked corpus. Not every sanction is scapegoating: discipline, restricted access, restitution, civil justice, and final judgment remain.

\Needspace{10\baselineskip}
\noindent\textbf{The Communion alternatives carry different clocks into new life.\enspace\elementref{Communion A; Communion B; Prayer after Communion}}
Communion A, if appointed, remembers wonders and food; Communion B hears Christ's present knock amid the Laodicean rebuke and anticipates table fellowship; the final prayer seeks new life. Compared as alternatives, memory makes past fidelity operative while interruption breaks complacency. Separate reading misses how each governs “new.” Its fruit is examination of what must be remembered and what summons is resisted. No exact conjunction was located. The branches are mutually alternative and unresolved; memory is not nostalgia, interruption is not rejection of tradition, and neither selection may be inferred.

\Needspace{10\baselineskip}
\noindent\textbf{Divine power refuses a single tempo.\enspace\elementref{First Reading; Second Reading; long Gospel}}
Wisdom locates mastery in leniency; the long Gospel spans delayed harvest, hidden growth, and leaven; Romans places aid within weakness and groaning. Together, power restrains, waits, grows, permeates, and intercedes. Separate reading may call patience inactive, growth merely gradual, or groaning deficient; the conjunction rejects speed and spectacle as measures. Its fruit is leadership distinguishing urgency from panic and quiet efficacy from passivity. Tertullian and official weeds preaching are near analogues, not this pattern. Mustard and leaven are long-form only; delay may be vicious, and charity may require immediate protection.
